\chapter{Boolean Analysis — Surface Examples}
%%% generated by the blueprint pipeline from TCSlib.BooleanAnalysis.GaussianSpace.SurfaceExamples
\section{Overview}

Declarations of the Lean module \texttt{TCSlib.BooleanAnalysis.GaussianSpace.SurfaceExamples}.

\section{Declarations}

\begin{theorem}[Surface area of disjoint intervals]
\lean{BooleanAnalysis.GaussianSpace.surfaceArea_disjoint_intervals}
\label{BooleanAnalysis.GaussianSpace.surfaceArea_disjoint_intervals}
\leanok
\uses{BooleanAnalysis.GaussianSpace.gaussianSurfaceArea, BooleanAnalysis.ThresholdFunctions.standardGaussianPDF}
Let $m\ge 0$, and let $a,b:\{1,\ldots,m\}\to\mathbb R$ satisfy $a_j<b_j$ for every $j$.
Suppose the closed intervals $[a_i,b_i]$ are pairwise disjoint. The Gaussian surface
area in $\mathbb R$ of their union is
$\sum_{j=1}^m\bigl(\varphi(a_j)+\varphi(b_j)\bigr)$, where $\varphi$ is the standard
Gaussian density.
\end{theorem}

\begin{theorem}[Gaussian surface area of a centred ball]
\lean{BooleanAnalysis.GaussianSpace.surfaceArea_centered_ball}
\label{BooleanAnalysis.GaussianSpace.surfaceArea_centered_ball}
\leanok
\uses{BooleanAnalysis.GaussianSpace.gaussianSurfaceArea}
Let $n\ge 1$ and $r>0$. The Gaussian surface area of the closed Euclidean ball of radius
$r$ centred at the origin in $\mathbb R^n$ is
\[\frac{2^{1-n/2}}{\Gamma(n/2)}\,r^{n-1}e^{-r^2/2}.\]
\end{theorem}

\begin{theorem}[Gaussian surface area bound for a ball]
\lean{BooleanAnalysis.GaussianSpace.surfaceArea_ball_le}
\label{BooleanAnalysis.GaussianSpace.surfaceArea_ball_le}
\leanok
\uses{BooleanAnalysis.GaussianSpace.gaussianSurfaceArea}
Let $n\ge 0$, let $c\in\mathbb R^n$, and let $r>0$. The Gaussian surface area of the
closed Euclidean ball of radius $r$ centred at $c$ is at most $\sqrt{2/\pi}$.
\end{theorem}

\begin{theorem}[Dimension-dependent bound for convex sets]
\lean{BooleanAnalysis.GaussianSpace.surfaceArea_convex_le}
\label{BooleanAnalysis.GaussianSpace.surfaceArea_convex_le}
\leanok
\uses{BooleanAnalysis.GaussianSpace.gaussianSurfaceArea}
There exists a constant $C>0$ such that, for every natural number $n$ and every
measurable convex set $A\subseteq\mathbb R^n$, the Gaussian surface area of $A$ is at
most
\[C n^{1/4}.\]
\end{theorem}

\begin{theorem}[Convex sets with large Gaussian surface area]
\lean{BooleanAnalysis.GaussianSpace.exists_convex_surfaceArea_ge}
\label{BooleanAnalysis.GaussianSpace.exists_convex_surfaceArea_ge}
\leanok
\uses{BooleanAnalysis.GaussianSpace.gaussianSurfaceArea}
There is a constant $c>0$ such that for every positive integer $n$ there exists a closed
convex set $A\subseteq\mathbb R^n$ whose Gaussian surface area is at least $c n^{1/4}$.
\end{theorem}

\begin{theorem}[Surface area of halfspace intersections]
\lean{BooleanAnalysis.GaussianSpace.surfaceArea_halfspace_intersection_le}
\label{BooleanAnalysis.GaussianSpace.surfaceArea_halfspace_intersection_le}
\leanok
\uses{BooleanAnalysis.GaussianSpace.gaussianSurfaceArea}
There exists a constant $C>0$ such that, for all natural numbers $n,m$ with $m\ge2$, all
vectors $a_1,\ldots,a_m\in\mathbb R^n$, and all thresholds $t_1,\ldots,t_m\in\mathbb R$,
the Gaussian surface area of $\{x\in\mathbb R^n: a_j\cdot x\le t_j\text{ for every }j\}$
is at most $C\sqrt{\log m}$.
\end{theorem}

\begin{theorem}[Halfspace intersections with large Gaussian surface area]
\lean{BooleanAnalysis.GaussianSpace.exists_halfspace_intersection_surfaceArea_ge}
\label{BooleanAnalysis.GaussianSpace.exists_halfspace_intersection_surfaceArea_ge}
\leanok
\uses{BooleanAnalysis.GaussianSpace.gaussianSurfaceArea}
There is a constant $c>0$ such that for every integer $m\ge2$, there are a dimension
$n\ge0$, vectors $a_1,\ldots,a_m\in\mathbb R^n$, and thresholds
$t_1,\ldots,t_m\in\mathbb R$ for which the intersection
\[A=\{x\in\mathbb R^n: a_j\cdot x\le t_j\text{ for every }j=1,\ldots,m\}\]
has Gaussian surface area at least $c\sqrt{\log m}$.
\end{theorem}

\begin{theorem}[Surface area bound for convex cones]
\lean{BooleanAnalysis.GaussianSpace.surfaceArea_convex_cone_le}
\label{BooleanAnalysis.GaussianSpace.surfaceArea_convex_cone_le}
\leanok
\uses{BooleanAnalysis.GaussianSpace.gaussianSurfaceArea}
Let $n\ge 0$ and let $A\subseteq\mathbb R^n$ be measurable and convex. Suppose that
whenever $x\in A$ and $t\ge 0$, the scalar multiple $tx$ also belongs to $A$. Then the
Gaussian surface area of $A$ is at most $1$.
\end{theorem}

\begin{theorem}[Surface area of polynomial threshold sets]
\lean{BooleanAnalysis.GaussianSpace.surfaceArea_polynomial_threshold_le}
\label{BooleanAnalysis.GaussianSpace.surfaceArea_polynomial_threshold_le}
\leanok
\uses{BooleanAnalysis.GaussianSpace.gaussianSurfaceArea}
Let $n,k\ge0$, and let $p$ be a real polynomial in $n$ variables whose total degree is
at most $k$. The Gaussian surface area of the set $\{x\in\mathbb R^n:p(x)\ge0\}$ is at
most $k/\sqrt{2\pi}$.
\end{theorem}

\begin{theorem}[Polynomial threshold sets with large Gaussian surface area]
\lean{BooleanAnalysis.GaussianSpace.exists_polynomial_threshold_surfaceArea_ge}
\label{BooleanAnalysis.GaussianSpace.exists_polynomial_threshold_surfaceArea_ge}
\leanok
\uses{BooleanAnalysis.GaussianSpace.gaussianSurfaceArea}
For every natural number $k$ and every real $\varepsilon>0$, there is a real polynomial
$p$ in one variable of degree at most $k$ such that the Gaussian surface area of the set
$\{x\in\mathbb R:p(x)\ge 0\}$ is at least
$\max\!\left(\frac{k}{\sqrt{2\pi}}-\varepsilon,0\right)$.
\end{theorem}
